\bookStart{Ribe galder stick (DR EM85;493)}
\setBookCode{RibeStick}

\begin{flushright}%
\textbf{Dating:} Mediæval.%TODO

\textbf{Meter:} \Fornyrdislag, \Galdralag%para
\end{flushright}%

\subsection{Introduction}

A wooden stick from the Danish city of Ribe.  The galder is syncretic and contains numerous pre-Christian elements in a Christian(ised) context.

The inscription may be conveniently divided into four parts.  Part one (ll.~1--4) contains an introductory prayer where the healer asks for the aid of natural forces (Earth, Up-heaven and the Sun) and Christian divinitities (God and Saint Mary) so that the healing may be successful.  Part two (ll.~5--8) ritually exorcises any sickness which may have entered any part of the body.  Part three (ll.~9--14) apparently warns the addressee that they will be haunted by “nine needs” (an old Heathen formula; see note) until they say the charm.  Part four (ll. 15, which is probably prose) gives the personal name “Bonde”, perhaps the addressee, and concludes with an “Amen”.

\subsection{Transliteration}

Note that the part of § C following ¶ is deliberately scraped off.

§ A \textbf{+ io\textasciicircum rþ ÷ biþ a\textasciicircum k ÷ ua\textasciicircum rþæ ÷ o\textasciicircum k ÷ uphimæn ÷ so\textasciicircum l ÷ o\textasciicircum k ÷ sa\textasciicircum nt\textasciicircum æ maria ÷ o\textasciicircum k ÷ salfæn ÷ gud| |drotæn ÷ þæt han ÷ læ mik ÷ læknæs÷ha\textasciicircum nd ÷ o\textasciicircum k lif÷tuggæ ÷ at\textasciicircum  \textasciicircum liuæ} \\
§ B \textbf{uiuindnæ ÷ þær ÷ botæ ÷ þa\textasciicircum rf ÷ or ÷ ba\textasciicircum k ÷ o\textasciicircum k or brʀst ÷ or lækæ ÷ o\textasciicircum k or lim ÷ or øuæn ÷ o\textasciicircum k or øræn ÷ or ÷ a\textasciicircum llæ þe ÷ þær ÷ ilt ÷ kan i at} \\
§ C \textbf{kumæ ÷ suart ÷ hetær ÷ sten ÷ ha\textasciicircum n ÷ stær ÷ i ÷ hafæ ÷ utæ ÷ þær ÷ ligær ÷ a ÷ þe ÷ ni ÷ no\textasciicircum uþær ÷ þæ¶r ÷ l-{}-{}-r(a) ÷ (þ)en-nþþæþeskulhuærki} \\
§ D \textbf{skulæ ÷ huærki ÷ søtæn ÷ sofæ ÷ æþ ÷ uarmnæn ÷ uakæ ÷ førr æn ÷ þu ÷ þæssa ÷ bot ÷ biþær ÷ þær ÷ a\textasciicircum k o\textasciicircum rþ ÷ at kæþæ ÷ ro\textasciicircum nti ÷ amæn ÷ o\textasciicircum k þæt ÷ se +}

\sectionline

\newpage

\subsection{Text}

\bvg\bva[]%
\edtrans{\alst{Jo}rð bið ak varðę \hld\ ok \alst{u}p-himęn}{I bid Earth to guard and Up-heaven}{\Bfootnote{\emph{Jorð \dots\ ok up-himęn} ‘Earth \dots\ and Up-heaven’ is a formulaic and originally pre-Christian merism representing the totality of the cosmos as separated from the primæval ettin-hermaphrodite Yimer; cf.~\textlink{Voluspa}[3].  It is here used in a syncretic magical context; the galer (galder-singer) asks the Cosmos itself to watch over and protect him/her as he/she utters the magic words.
A notable parallel to the present instance is found in the OE \textlink{Acerbot}[1]/4 (ca.~1000 \ce), edited above.  Cf.~also a 14th c.~version of the Icelandic \emph{Griða-mál} ‘Grith-speech’, a legal formula for the establishments of truces between former enemies, where we read that: \emph{Jǫrð rę́ðr griðum fyr neðan, en upp-himinn varðar fyr ofan, en sjór fyr útan, sá er kringir um ǫll lǫnd.} ‘The earth rules the griths from below, and Upheaven guards them from above, and the sea from without which encircles all lands.’ \parenciteshorthand[659]{DI_2}.
Due to its invocation of pre-Christian cosmology along with later use of pre-Christian legal formulae like \emph{rę́kr ok rekinn} and \emph{grid-nídings nafn} it is likely that the \emph{Grida-mál} reflects a late remnant of pre-Christian law, and together with the present galder and \textlink{Acreboot}[1]/4 all three texts appear to reflect a common pre-Christian Germanic ritual expression.}} &
\alst{s}ól ok \alst{s}antę María \hld\ ok \alst{s}alfęn Guð dróttęn &
þęt hann \alst{l}ę́ mik \alst{l}ę́knęs-hand \hld\ ok \alst{l}yf-tungę &
at lyfę \emph{\alst{b}}ifj\emph{a}ndę \hld\ þęr \alst{b}ótę þarf.\eva

\bvb I bid Earth to guard and Up-heaven, \\
the Sun and Saint Mary, and the very Lord God \\
that he lend me a leecher’s hand and a medicine-tongue \\
as medicine for the trembler who needs the cure.\evb\evg


\bvg\bva[][5]%
\ind Ór \alst{b}ak ok ór \alst{b}rýst &
\ind ór \alst{l}íkę ok ór \alst{l}im &
\ind ór \alst{ø̂}vęn ok ór \alst{ø̂}ręn &
ór \alst{a}llę þé \hld\ þęr \alst{i}llt kann í at·kumę.\eva

\bvb Out of back and out of breast! \\
Out of body and out of limb! \\
Out of eyes and out of ears! \\
Out of everything, where evil might come in!\evb\evg


\bvg\bva[][9]%
Svart hêtęr \alst{st}ênn \hld\ hann \alst{st}ę́r í hafę útę, &
\ind þęr liggęr á þé \alst{n}í\emph{u} \alst{n}auðęr; &
\ind \edtext{þę́}{\Afootnote{add., but scraped off and poorly legible \emph{\textbf{†÷ þæ¶r ÷ l-{}-{}-r(a) ÷ (þ)en-nþþæþeskulhuærki†}} (a poorly carved form of ll.~9--10) \emph{DR EM85;493}}} skulę hvęrki \alst{s}ǿtęn \alst{s}ofę; &
\ind ęð \alst{v}armęn \alst{v}akę; &
\ind førr ęn þú þęssa \alst{b}ót \alst{b}iðęr, &
þęr \alst{a}k \alst{o}rð \hld\ \alst{a}t·k\emph{v}ę́ðę r\emph{ú}nti. &
\ind Amęn ok þat sé.\eva

\bvb Swart is a stone called; it stands out in the ocean. \\
\ind There lie on it nine needs; \\
\ind they will neither sleep sweetly \\
\ind nor wake warmly, \\
\ind until thou prayest this cure \\
to which I have spoken the runic(?) words. \\
\ind Amen and let it be so.\evb\evg

\sectionline
